% ============================================
% DOC-12 (EN)
% English translation — SPEC-935-R355
% ============================================
\section*{\textbf{DOC-12} --- Oliveira's Diary II}
\markboth{DOC-12 --- Oliveira's Diary II}{}

\noindent \textit{The handwriting changes}

\noindent \textbf{Dr. Heitor Oliveira --- Personal Notebook}
\textbf{April 20, 1979}

\vspace{0.5em}\noindent\rule{\textwidth}{0.5pt}\vspace{0.5em}

\noindent Day 7.

\noindent My handwriting is changing.

\noindent I noticed this morning, when I was taking notes about a patient. I looked at the chart and did not recognize the handwriting. It seemed more... crooked. Childish. Like the handwriting of someone who learned to write on his own.

\noindent Like Joaquim's handwriting.

\noindent I took his diary --- yes, it is still with me, I cannot get rid of it --- and compared. The similarity is striking. The capital \textquotedblleft{}J\textquotedblright{}, the lowercase \textquotedblleft{}a\textquotedblright{}, the final stroke of the words.

\noindent I am copying his handwriting without noticing.

\noindent Or he is copying mine.

\noindent Or the same hand is writing both.

\vspace{0.5em}\noindent\rule{\textwidth}{0.5pt}\vspace{0.5em}

\noindent Day 8.

\noindent Today I tried to get rid of the diary.

\noindent I put it inside a plastic bag. I tied it. I drove to a bus stop on the BR-040, 15 kilometers from Barbacena. I threw the bag into a public trash bin.

\noindent I went home.

\noindent The diary was on the kitchen table.

\noindent No bag. No ties. Just the diary, open at the page where I had stopped reading.

\noindent My wife said: \textquotedblleft{}Did someone bring it?\textquotedblright{}

\noindent I said: \textquotedblleft{}No.\textquotedblright{}

\noindent She said: \textquotedblleft{}Then how did it come back?\textquotedblright{}

\noindent I did not answer.

\noindent I picked up the diary. I put it in my coat pocket. This time, I will not try to get rid of it. There is no point.

\noindent It wants to be read.
It \textit{needs} to be read.
And I need to read.

\vspace{0.5em}\noindent\rule{\textwidth}{0.5pt}\vspace{0.5em}

\noindent Day 9.

\noindent I read 50 more pages.

\noindent Joaquim describes the encounter at the trench in details that no delusion could produce. The creature spoke to him. Made a deal. Gave the hunger in exchange for company.

\noindent There is a passage that chilled me:

\noindent \textit{\textquotedblleft{}She said: 'One day, a man in a white coat will read this diary. He will deny it. He will be afraid. He will try to run away. But he will come back. Because curiosity is greater than fear. And when he comes back, I will be waiting.'\textquotedblright{}}

\noindent \textit{\textquotedblleft{}I asked: 'How do you know?'\textquotedblright{}}

\noindent \textit{\textquotedblleft{}She said: 'Because he has been here before. In another life. In another body. He always comes back. From the beginning.'\textquotedblright{}}

\noindent From the beginning.

\noindent How many readers came before me?

\noindent How many doctors, researchers, curious souls opened this diary and felt the smell, and saw the eye, and began to dream?

\noindent Where are they now?

\vspace{0.5em}\noindent\rule{\textwidth}{0.5pt}\vspace{0.5em}

\noindent Day 10.

\noindent I did research.

\noindent The diary was mentioned in three academic articles in the 1980s. All the authors had \textquotedblleft{}health problems\textquotedblright{} before publishing. One died of cancer. Another had a nervous breakdown. The third --- a psychiatrist named Dr. Verônica Tavares --- disappeared in 1985.

\noindent The diary was found in her house. Together with a note:

\noindent \textit{\textquotedblleft{}The Cycle cannot be broken. Only passed on. I return it to the archive so that another may read. I am sorry.\textquotedblright{}}

\noindent She passed it on.

\noindent And I took it.

\vspace{0.5em}\noindent\rule{\textwidth}{0.5pt}\vspace{0.5em}

\noindent Day 11.

\noindent 3:14 in the morning.

\noindent I woke up with my eye wide open.

\noindent That is not a metaphor. The right eye was open. I could not close it. The eyelid was stiff.

\noindent I went to the mirror.

\noindent The right eye is yellow.

\noindent I will not go to a doctor. I will not tell anyone. I have read enough to know what is happening.

\noindent The creature is passing into me.

\noindent Joaquim died.
I read the diary.
Now she is in me.

\noindent My eye is yellow.
My smell is sweet.
My handwriting is like his.

\noindent I am becoming patient 1,261.

\vspace{0.5em}\noindent\rule{\textwidth}{0.5pt}\vspace{0.5em}

\noindent Day 12.

\noindent I have decided.

\noindent I will write my own diary. I will record what is happening to me. So that, when someone reads it, they will know.

\noindent The Cycle is a chain of readers.

\noindent Joaquim → Oliveira → ? → ? → ?

\noindent Someone will read this after me.

\noindent If you are reading: I am sorry.

\noindent There is no way to stop it.

\noindent The diagnosis is not a metaphor.

\noindent The creature is real.

\noindent I saw the eye.

\noindent I felt the smell.

\noindent I am the next one.

\noindent \textit{--- Heitor Oliveira, April 29, 1979}
\textit{--- Patient 1,261}

\vspace{0.5em}
\noindent\textit{\textbf{↪ Routes:} \rota{DOC-13} \rotasep \rota{CONT-01}}
